% TODO make hw stylesheet
\documentclass[11pt]{scrartcl}
\usepackage[a4paper, total={6in, 8in}]{geometry}
\usepackage[usenames,dvipsnames,svgnames]{xcolor}
\usepackage[shortlabels]{enumitem}
\usepackage[framemethod=TikZ]{mdframed}
\usepackage{amsmath,amssymb,amsthm}
\usepackage[colorlinks]{hyperref}
\usepackage{multicol}
\usepackage{tabularx}

\hypersetup{
colorlinks=true,
linkcolor=blue,
filecolor=magenta,      
urlcolor=magenta,
pdftitle={MAT 319 HW}
}

\usepackage{microtype}
\usepackage{mathtools}
\usepackage[headsepline]{scrlayer-scrpage}
\usepackage{thmtools}
\usepackage{todonotes}

\mdfdefinestyle{mdgreenbox}{linecolor=ForestGreen,backgroundcolor=ForestGreen!5,
	linewidth=2pt,rightline=false,leftline=true,topline=false,bottomline=false,}
\declaretheoremstyle[headfont=\bfseries\sffamily\color{ForestGreen!70!black},
mdframed={style=mdgreenbox},headpunct={ --- },]{thmgreenbox}

\declaretheorem[style=thmgreenbox,name=Claim,numbered=no]{claim*}
\declaretheorem[style=thmgreenbox,name=Lemma,numbered=no]{lemma*}

\addtolength{\textheight}{3.14cm}
\providecommand{\ol}{\overline}
\providecommand{\eps}{\varepsilon}
\providecommand{\half}{\frac{1}{2}}
\providecommand{\dang}{\measuredangle} %% Directed angle
\providecommand{\CC}{\mathbb C}
\providecommand{\FF}{\mathbb F}
\providecommand{\NN}{\mathbb N}
\providecommand{\QQ}{\mathbb Q}
\providecommand{\RR}{\mathbb R}
\providecommand{\ZZ}{\mathbb Z}
\providecommand{\dg}{^\circ}
\providecommand{\ii}{\item}
\providecommand{\alert}[1]{{\sffamily\textbf{\textcolor{blue}{#1}}}}

\reversemarginpar

\theoremstyle{problem}
\newtheorem{innercustomthm}{Problem}
\newenvironment{problem}[1]
{\renewcommand\theinnercustomthm{#1}\innercustomthm}
{\endinnercustomthm}

\theoremstyle{remark}
\newtheorem{remarknumbered}{Remark}
\newenvironment{remark}
{\renewcommand\theremarknumbered{\hspace{-\the\fontdimen2\font\space}}\remarknumbered}
{\endremarknumbered}

\newenvironment{soln}{\begin{proof}[Solution]}{\end{proof}}
\newenvironment{walkthrough}{\noindent \textbf{Walkthrough.}}{\medskip}
\newlist{walk}{enumerate}{3}
\setlist[walk]{label=\bfseries (\alph*)}

\providecommand{\pow}{\operatorname{Pow}}
\providecommand{\vocab}[1]{{\textbf{\textcolor{ForestGreen}{#1}}}}
\providecommand{\lcm}{\operatorname{lcm}}
\providecommand{\abs}[1]{\left|#1\right|}

\usepackage{asymptote}
\begin{asydef}
defaultpen(fontsize(10pt));
unitsize(1cm);
size(8cm); // set a reasonable default
usepackage("amsmath");
usepackage("amssymb");
settings.tex="pdflatex";
settings.outformat="pdf";
// Replacement for olympiad+cse5 which is not standard
import geometry;
// recalibrate fill and filldraw for conics
void filldraw(picture pic = currentpicture, conic g, pen fillpen=defaultpen, pen drawpen=defaultpen)
{ filldraw(pic, (path) g, fillpen, drawpen); }
void fill(picture pic = currentpicture, conic g, pen p=defaultpen)
{ filldraw(pic, (path) g, p); }
// some geometry
pair foot(pair P, pair A, pair B) { return foot(triangle(A,B,P).VC); }
pair orthocenter(pair A, pair B, pair C) { return orthocentercenter(A,B,C); }
pair centroid(pair A, pair B, pair C) { return (A+B+C)/3; }
// cse5 abbreviations
path CP(pair P, pair A) { return circle(P, abs(A-P)); }
path CR(pair P, real r) { return circle(P, r); }
pair IP(path p, path q) { return intersectionpoints(p,q)[0]; }
pair OP(path p, path q) { return intersectionpoints(p,q)[1]; }
path Line(pair A, pair B, real a=0.6, real b=a) { return (a*(A-B)+A)--(b*(B-A)+B); }
// cse5 more useful functions
picture CC() {
	picture p=rotate(0)*currentpicture;
	currentpicture.erase();
	return p;
}
pair MP(Label s, pair A, pair B = plain.S, pen p = defaultpen) {
	Label L = s;
	L.s = "$"+s.s+"$";
	label(L, A, B, p);
	return A;
}
pair Drawing(Label s = "", pair A, pair B = plain.S, pen p = defaultpen) {
	dot(MP(s, A, B, p), p);
	return A;
}
path Drawing(path g, pen p = defaultpen, arrowbar ar = None) {
	draw(g, p, ar);
	return g;
}
\end{asydef}

\usepackage{mathrsfs}

\ihead{\footnotesize MAT 319 HW03}
\ohead{\footnotesize September 18, 2024}

\title{\vspace{-2em}MAT 319 HW03}
\author{\large Aditya Pahuja}
\date{\large Due 9/18}
\begin{document}
\maketitle

\begin{problem}{10.1}
	Which of the following sequences are increasing? decreasing? bounded?
	\begin{enumerate}[\emph{\textbf{(\alph*)}}]
		\ii $1/n$
		\ii $(-1)^n/n$
		\ii $n^5$
		\ii $\sin(n\pi/7)$
		\ii $(-2)^n$
		\ii $n/3^n$
	\end{enumerate}
\end{problem}
\begin{soln}
	(c) is increasing; (a) and (f) are decreasing; (a), (b), (d), and (f)
	are bounded.
	
	We deal with the monotonicity of the sequences first.
	
	First, (c) is increasing because
	\[(n + 1)^5 = n^5 + 5n^4 + 10n^3 + 10n^2 + 5n + 1 > n^5\]
	for all natural $n$.
	
	Second, (a) is decreasing because
	\[\frac{1}{n} > \frac{1}{n + 1}\]
	for all natural $n$ (since the denominator on the left side
	is bigger than that on the right), and (f) is decreasing because
	\[\frac{n}{3^n} = \frac{3n}{3^{n+1}} > \frac{n + 1}{3^{n+1}},\]
	where the inequality holds due to $3n > n + 1 \iff 2n > 1$
	for all natural $n$.
	
	Third, (b), (d), and (e) are not monotonic.
	In general, to show that a sequence $(s_n)$ isn't monotonic,
	it suffices to find pairs of elements $(s_a, s_{a+1})$
	and $(s_b, s_{b+1})$ such that
	\[s_a < s_{a+1},\qquad s_b > s_{b+1},\]
	since the first inequality shows that $s_n$ cannot be decreasing
	and the second shows that it cannot be increasing.
	Noting this, (b) isn't monotonic because
	\[-1 = \frac{(-1)^1}{1} < \frac{(-1)^2}{2} = \frac12,\qquad
	\frac12 = \frac{(-1)^2}{2} > \frac{(-1)^3}{3} = -\frac13\]
	(so $a = 1$ and $b = 2$); (d) isn't monotonic because
	\[\sin(6\pi/7) > \sin(7\pi/7) = 0,\qquad
	0 = \sin(14\pi/7) < \sin(15\pi/7),\]
	since $\sin(x)$ is positive whenever it can be written as
	$2\pi k + r$ where $k\in\ZZ$ and $r\in(0,\pi)$;
	and (f) isn't monotonic because
	\[-2 = (-2)^1 < (-2)^2 = 4,\qquad
	4 = (-2)^2 > (-2)^3 = -8.\]
	
	Now, we determine which sequences are bounded.
	
	First, (a) is bounded because $0 < 1/n \le 1$ for all $n$;
	(b) is bounded because $-1 \le (-1)^n/n \le 1$ for all $n$ (equivalently,
	$|(-1)^n/n| = 1/n \le 1$ for all $n$);
	(d) is bounded because $-1 \le \sin x < 1$ for all real $x$;
	and (f) is bounded because $0 < n/3^n < 1$, since $3^n > n$ for all $n$
	(this is easy to show by induction: it holds trivially for $n=1$,
	and $3^n > n$ implies $3^{n+1} = 3\cdot 3^n > 3n > n$ for all natural $n$).
	
	It remains to show that (c) and (e) are unbounded.
	This follows easily from noting that $n^5$ and $|(-2)^n| = 2^n$
	both diverge to $+\infty$.
\end{soln}

\pagebreak
\begin{problem}{10.2}
	Prove Theorem 10.2 for bounded decreasing sequences.
\end{problem}
\begin{soln}
	Let $(s_n)$ be a bounded decreasing sequence. We show that
	\[\lim_{n\to\infty} s_n = s\]
	where $s = \inf\{s_n\}$ is the infimum of the sequence $(s_n)$.
	In particular, we want to show that for all $\eps > 0$,
	there exists an $N$ such that $|s - s_n| < \eps$ for all $n > N$.
	
	First, there exists some $m$ such that $|s - s_m| < \eps$;
	this is because, if the above statement was false
	(i.e. $|s - s_m| \ge \eps$ for all $m$), then we can find a greater
	upper bound for $\{s_n\}$ in the form of $s + \eps/2$, which would
	contradict the fact that $s$ is the greatest upper bound of $\{s_n\}$.
	
	Then, since $s_n$ is decreasing, if $k > m$, then
	$s \le s_k \le s_m$, so
	\[\eps = |s_m - s| > |s_k - s| = |s - s_k|\]
	for all $k \ge m$, meaning we can take $N = m$ in the first paragraph.
	Thus, we have shown that $s_n$ converges to $s$, since this bounding
	works for any arbitrary $\eps$.
\end{soln}
\pagebreak
\begin{problem}{10.7}
	Let $S$ be a bounded nonempty subset of $\RR$ such that
	$\sup S$ is not in $S$. Prove that there is a sequence $(s_n)$
	of points in $S$ such that $\lim s_n = \sup S$.
\end{problem}

\begin{proof}
	Define $s = \sup S$.
	Let $s_1$ be an arbitrary element of $S$; we will then
	inductively define the rest of the sequence.
	
	Given $s_1$, \dots, $s_n$, define
	\[d_n = \min\{s - s_1, s - s_2, s - s_3, \dots, s - s_n\}\]
	(of course, each $s - s_i$ is positive because $s > x$ for all $x\in S$).
	Then, let $s_{n+1}$ be an element of $S$
	such that $s - s_{n+1} < \frac{d_n}{2}$; such an $s_{n+1}$ exists
	because $s$ is the supremum of $S$ (otherwise we would be able to
	find a smaller upper bound for $S$, such as $s - \frac{d_n}{4}$).
	In fact, this implies that $d_k = s - s_k$ for all natural $k$,
	since $s_{n + 1}$ is chosen specifically so that
	$s - s_{n + 1} < s - s_i$ for all $i\in\{1, 2, \dots, n\}$.
	
	I then claim that $\lim_{n\to\infty} s_n = s$.
	To show this, we want to prove that for each $\eps > 0$,
	$|s - s_n| = d_n < \eps$ eventually.
	First, we bound $d_n$ from above:
	
	\begin{claim*}
		For all $n > 1$,
		\[d_n < \frac{d_1}{2^{n - 1}}.\]
	\end{claim*}
	\begin{proof}
		We prove this by induction on $n$, noting in general that
		$d_{n+1} < \frac{d_n}{2}$ by definition
		(recall that $s_{n+1}$ is defined so that
		$s - s_{n+1} < \frac{d_n}{2}$).
		
		The base case $n = 2$ is obvious, since
		\[d_2 > \frac{d_1}{2} = \frac{d_1}{2^{2-1}}.\]
		
		Then, if $d_n < \frac{d_1}{2^{n - 1}}$,
		\[d_{n+1} < \frac{d_n}{2} < \frac{d_1}{2^{n-1}}\cdot\half
		= \frac{d_1}{2^n},\]
		so the claim is proven by the principle of mathematical induction.
	\end{proof}
	
	This means that it suffices to show that
	\[\frac{d_1}{2^{n-1}} < \eps\]
	is eventually true because the left-hand side is bigger than $d_n$.
	We can deduce this from the fact that, if $|a| < 1$,
	then $\lim_{n\to\infty} a^n = 0$, since this tells us that for all
	$\eps' > 0$, $|a^n - 0| = |a^n|$ is eventually less than $\eps'$.
	Picking $a = \half$ and $\eps' = \frac{\eps}{2d_1}$ shows us that
	\[\frac{1}{2^n} < \frac{\eps}{2d_1}
	\iff \frac{d_1}{2^{n-1}} < \eps\]
	eventually.
	
	Therefore, the limit of our $(s_n)$ sequence is $s$, as desired.
\end{proof}

\begin{remark}
	Since $(s_n)$ is monotonically increasing and bounded
	by construction, the proof that $\lim s_n = s$ is really just
	a proof that $s = \sup\{s_n\}$, since we are just showing that
	there exist elements of $\{s_n\}$ that are arbitrarily close to $s$
	(so $s$ is the least upper bound).
\end{remark}

\pagebreak
\begin{problem}{10.10}
	Let $s_1 = 1$ and $s_{n+1} = \frac13 (s_n + 1)$ for $n\ge 1$.
	\begin{enumerate}[\emph{\textbf{(\alph*)}}]
		\ii Find $s_2$, $s_3$, and $s_4$.
		\ii Use induction to show that $s_n > \half$ for all $n$.
		\ii Show that $(s_n)$ is a decreasing sequence.
		\ii Show $\lim s_n$ exists and find $\lim s_n$.
	\end{enumerate}
\end{problem}

\begin{soln}
	For (a),
	\[s_2 = \frac13(1 + 1) = \frac23,\quad
	s_3 = \frac13\left(\frac23+1\right) = \frac59,\quad
	s_4 = \frac13\left(\frac59 + 1\right) = \frac{14}{27}.\]
	For (b), we see that, clearly, $s_1 = 1$, and, if $s_n > \half$, then
	\[s_{n+1} = \frac13(s_n + 1) > \frac13\left(\half + 1\right) = \frac12,\]
	so, by the principle of mathematical induction, $s_n > \half$ for all $n$.
	
	To do parts (c) and (d), we will prove the following statement:
	\begin{claim*}
		For all $n > 1$,
		\[s_n = \frac{1}{2\cdot 3^{n-1}} + \half.\]
	\end{claim*}
	\begin{proof}
		Rewrite the recurrence relation as
		\[3s_n = s_{n-1} + 1 \iff 3\left(s_n - \half\right) = s_{n-1} - \half\]
		where the second equation is a result of subtracting $\frac32$
		from both sides.
		Then, if $(t_n)$ is the sequence given by $t_n = s_n - \half$
		for all $n\in\NN$, we can write
		\[3t_n = t_{n-1}\]
		for all $n > 1$.
		Then, by induction, we can say that $t_n = \frac{t_1}{3^{n-1}}$:
		the base case $n = 1$ is obvious ($t_1 = \frac{t_1}{3^0}$),
		and, for the inductive step, if $t_n = \frac{t_1}{3^{n-1}}$, then
		\[t_{n+1} = \frac{t_n}{3} = \frac{t_1}{3^{n-1}\cdot 3}
		= \frac{t_1}{3^n}.\]
		
		Thus, we can say that, for all $n$,
		\[s_n - \half = t_n = \frac{t_1}{3^{n-1}},\]
		and we know that $t_1 = s_1 - \half = \half$, which gives us
		\[s_n = \frac{1}{2\cdot 3^{n-1}} + \half.\]
	\end{proof}

	Now, to show (c) (that $(s_n)$ is decreasing), we simply observe that
	\[s_{n} - s_{n+1} =
	\left(\frac{1}{2\cdot 3^{n-1}} + \half\right) -
	\left(\frac{1}{2\cdot 3^n} + \half\right) =
	\frac{3}{2\cdot 3^{n}} - \frac{1}{2\cdot 3^{n}} =
	\frac{1}{3^{n}} > 0,\]
	implying $s_{n+1} < s_n$ for all $n \ge 2$ (and obviously
	$\frac23 = s_2 < s_1 = 1$ is true).
	
	Finally, for part (d),
	we can see very easily that the limit of $s_n$ is $\half$,
	since, as $n$ goes to infinity, the limit of $\frac{1}{3^{n-1}}$ is zero.\\
	
	Also, here's another (slightly quicker?) solution path for (c) and (d)
	based on the problem shown in recitation.
	For (c), induct on $n$ to show $s_n > s_{n+1}$ for all $n$,
	with $n=1$ giving the obvious $1 = s_1 > s_2 = \frac23$
	and $s_n > s_{n+1}$ implying
	\[s_{n+1} = \frac{s_n + 1}{3} > \frac{s_{n+1} + 1}{3} = s_{n+2}.\]
	For (d), we see that the limit of $s_n$ exists because $s_n$
	is bounded (each $s_n$ is in the interval $[\half, 1]$ by the work above)
	and monotonically decreasing, so, letting $L = \lim s_n$, we get
	\[L = \lim_{n\to\infty} s_n = \lim_{n\to\infty}\frac{s_{n-1} + 1}{3}
	= \frac{\lim_{n\to\infty} s_{n-1} + \lim_{n\to\infty} 1}
	{\lim_{n\to\infty}3} = \frac{L + 1}{3},\]
	which can easily be solved to get $L = \half$.
\end{soln}

\pagebreak
\begin{problem}{11.2, 11.3}
	Consider the following eight sequences:
	\[a_n = (-1)^n,\quad b_n = \frac 1n,\quad c_n = n^2,\quad
	d_n = \frac{6n + 4}{7n - 3},\]
	\[s_n = \cos\left(\frac{\pi n}{3}\right),\quad
	t_n = \frac{3}{4n + 1},\quad u_n = \left(-\half\right)^n,
	\quad v_n = (-1)^n + \frac 1n.\]
	\begin{enumerate}[\emph{\textbf{(\alph*)}}]
		\ii For each sequence, give an example of a monotone subsequence.
		\ii For each sequence, give its set of subsequential limits.
		\ii For each sequence, give its $\limsup$ and $\liminf$.
		\ii Which of the sequences converges? diverges to $+\infty$?
		diverges to $-\infty$?
		\ii Which of the sequences is bounded?
	\end{enumerate}
\end{problem}

\begin{soln}
\begin{enumerate}[\textbf{(\alph*)}]
	\ii $(b_n)$, $(c_n)$, $(d_n)$, and $(t_n)$ are all already monotonic,
	so the sequences themselves (or any subsequence at all) suffice as answers.
	The only nontrivial one of these to verify is $d_n$;
	for completeness, we verify this by noting that
	\[d_n = \frac{6}{7} + \frac{10}{49n-21},\]
	which is decreasing because $49n-21$ is always positive and increasing.
	$(a_n)$ and $(s_n)$ are periodic, so for example
	\[(1, 1, 1, 1, \dots) = (a_2, a_4, a_6, a_8, \dots) =
	(s_6, s_{12}, s_{18}, s_{24}, \dots)\]
	is an example of a monotone subsequence,
	and finally, for $u_n$, we can take
	\[\left(\frac14, \frac1{16}, \frac1{64}, \frac1{256}, \dots\right)
	= (u_2, u_4, u_6, u_8, \dots)\]
	and for $v_n$ we can take
	\[\left(\frac{3}{2}, \frac{5}{4}, \frac76, \frac98, \dots\right)
	= (v_2, v_4, v_6, v_8, \dots).\]
	\ii Big list:
	\begin{itemize}
		\ii The set of subsequential limits of $(a_n)$ is $\{-1, 1\}$.
		\ii The set of subsequential limits of $(b_n)$ is $\{0\}$.
		\ii The set of subsequential limits of $(c_n)$ is $\{+\infty\}$.
		\ii The set of subsequential limits of $(d_n)$ is $\{\frac67\}$.
		\ii The set of subsequential limits of $(s_n)$ is
		$\{-1, -\frac12, \frac12, 1\}$.
		\ii The set of subsequential limits of $(t_n)$ is $\{0\}$.
		\ii The set of subsequential limits of $(u_n)$ is $\{0\}$.
		\ii The set of subsequential limits of $(v_n)$ is $\{-1, 1\}$.
	\end{itemize}
	\ii We can just look at the previous part's answers and compute
	the supremum/infimum of each set of subsequential limits:
	\begin{itemize}
		\ii $\limsup a_n = 1$ and $\liminf a_n = -1$
		\ii $\limsup b_n = 0$ and $\liminf b_n = 0$
		\ii $\limsup c_n = +\infty$ and $\liminf c_n = +\infty$
		\ii $\limsup d_n = \frac67$ and $\liminf d_n = \frac67$
		\ii $\limsup s_n = 1$ and $\liminf s_n = -1$
		\ii $\limsup t_n = 0$ and $\liminf t_n = 0$
		\ii $\limsup u_n = 0$ and $\liminf u_n = 0$
		\ii $\limsup v_n = 1$ and $\liminf v_n = -1$
	\end{itemize}
	\ii Convergence/divergence to $\pm\infty$
	can be seen by whether the $\limsup$ and $\liminf$ are the same value.
	\begin{itemize}
		\ii $(a_n)$ neither converges nor diverges.
		\ii $(b_n)$ converges to $0$.
		\ii $(c_n)$ diverges to $+\infty$.
		\ii $(d_n)$ converges to $\frac67$.
		\ii $(s_n)$ neither converges nor diverges.
		\ii $(t_n)$ converges to $0$.
		\ii $(u_n)$ converges to $0$.
		\ii $(v_n)$ neither converges nor diverges.
	\end{itemize}
	\ii For any sequence $(x_n)$,
	if both $\limsup x_n$ and $\liminf x_n$ are real numbers,
	then $(x_n)$ is bounded.
	Thus, all the sequences except $(c_n)$ are bounded.
\end{enumerate}
\end{soln}

\pagebreak
\begin{problem}{11.5}
	Let $(q_n)$ be an enumeration of all the rationals in the interval $(0, 1]$.
	\begin{enumerate}[\emph{\textbf{(\alph*)}}]
		\ii Give the set of subsequential limits for $(q_n)$.
		\ii Give the values of $\limsup q_n$ and $\liminf q_n$.
	\end{enumerate}
\end{problem}

\begin{soln}
	The answer to (a) is the interval \framebox{$[0,1]$}.
	Let $S$ be the set of subsequential limits.
	
	We start by showing that $S\subseteq[0,1]$.
	In particular, for any sequence $(x_i)$ with $0 < x_i \le 1$ for all $x_i$,
	we show that, if $\lim x_i$ exists, then it must be in $[0, 1]$.
	Suppose otherwise, that is, the limit $L$ is either less than zero
	or greater than $1$. If $L > 1$, then $\lim_{n\to\infty} x_n = L$
	means that there exists an $x_i$ such that
	\[L - x_i = |L - x_i| < L - 1\]
	(where we use the fact that $x_i \le 1 < L$ in the equality);
	this gives the contradiction $-x_i < -1 \iff 1 < x_i$.
	Similarly, if $L < 0$, then there exists an $x_i$ such that
	\[x_i - L = |x_i - L| < 0 - L\]
	(using the fact that $L < 0 < x_i$),
	again giving a contradiction in the form $x_i < 0$.
	
	Now, we show that $[0, 1]\subseteq S$, i.e. any $r\in[0, 1]$
	is a subsequential limit.
	We will make use of the following lemma,
	which was on homework 1:
	\begin{lemma*}[Exercise 4.11 in the textbook]
		For any real $a$, $b$ with $a < b$,
		there are infinitely many rationals between $a$ and $b$.
	\end{lemma*}
	\begin{proof}
		Suppose that there are finitely many rationals in $(a, b)$.
		Then, let $q$ be the largest one. Clearly, the interval $(q, b)$
		cannot have any rational numbers in it --- however,
		this contradicts the density of $\QQ$ in $\RR$, contradiction.
	\end{proof}

	Our goal is to construct a sequence of indices $(n_k)$
	such that the subsequence of $(q_n)$ given by $(q_{n_k})_{k\in\NN}$
	converges to $r$. First, set $n_1 = 1$.
	Then, to get $n_{k+1}$, we will find a rational number
	that appears after $q_{n_k}$ in the $(q_n)$ sequence (i.e. $n_{k+1} > n_k$)
	such that $|r - q_{n_{k+1}}| < \frac{1}{2^{k+1}}$.
	
	To show that such a rational number exists, we note that the interval
	\[\left(r - \frac{1}{2^{n+1}}, r + \frac{1}{2^{n+1}}\right)
	\cap[0, 1]\]
	contains infinitely rational numbers between $0$ and $1$.
	Since there are only finitely rational numbers
	appearing at or before $q_{n_k}$ in the $(q_n)$ sequence
	(namely $q_1$, \dots, $q_{n_k}$), there must exist a rational number
	$q\in I$ such that the natural number $m$ for which $q_m = q$
	is greater than $n_k$; set $n_{k+1} = m$.
	
	In this way, we constructed an increasing sequence $(n_k)$ of naturals,
	which induces the subsequence $(q_{n_k})$ of $(q_n)$.
	Moreover, we know that, for each $\eps > 0$, there exists a $K$
	such that $\eps > \frac{1}{2^k}$ for all $k > K$, so indeed
	\[\abs{q_{n_k} - r} < \frac{1}{2^k} < \eps\]
	for $k > K$, hence the limit of the $(q_{n_k})$ subsequence
	is indeed $r$, as desired.
	
	Since $S$ and $[0, 1]$ contain each other, they are the same set,
	which is what we claimed.
	
	Part (b) is then obvious, since $\limsup q_n$ and $\liminf q_n$
	are respectively $\sup[0, 1] = 1$ and $\inf[0, 1] = 0$.
\end{soln}

\pagebreak
\begin{problem}{11.8}
	Prove that $\liminf s_n = -\limsup (-s_n)$ for any sequence $(s_n)$.
\end{problem}
\begin{soln}
	We use this lemma relating supremums and infimums
	(suprema and infima?):
	\begin{lemma*}[Exercise 5.4 in the textbook]
		For any set $S$, $\inf S = -\sup (-S)$, where $-S$ is the set
		$\{-s : s\in S\}$.
	\end{lemma*}
	\begin{proof}
		By definition, $c = \inf S$ is the greatest lower bound of $S$.
		Since $c \le x$ for all $x\in S$, $-c \ge -x$ for all $x\in S$,
		or equivalently $-c \ge y$ for all $y\in -S$, so $-c$ is an upper bound
		for $-S$, in particular $-c\ge \sup(-S)$.
		
		Suppose for contradiction that $-c > \sup(-S)$.
		Then, there exists an $\eps$ for which
		every element of $-S$ is at least $\eps$ away from $-c$;
		that is, $|-c - y| \ge \eps$ for all $y\in -S$.
		However, this means that $|(-y) - c| \ge \eps$ for all $y\in -S$,
		or equivalently
		\begin{equation}\label{11-8:bigcontradiction}
			|x - c| \ge \eps \text{ for all } x\in S.
		\end{equation}
		
		Now, since $c = \inf S$, we know that, for every $\hat\eps$,
		there exists an $x_{\hat\eps}\in S$ such that
		\[|x_{\hat\eps} - c| < \hat\eps.\]
		This gives the desired contradiction by taking $\hat\eps = \eps$
		and $x = x_\eps$ in (\ref{11-8:bigcontradiction}).
	\end{proof}

	To that end, we have
	\[\liminf s_n = \lim_{n\to\infty} \inf\{s_k : k\ge n\}
	= \lim_{n\to\infty} -\sup\{-s_k : k\ge n\}
	= -\lim_{n\to\infty} \sup\{-s_k : k\ge n\}.\]
	The latter expression is then $-\limsup (-s_n)$,
	so $\liminf s_n = -\limsup (-s_n)$ as desired. 
\end{soln}

\end{document}